\section{Numerical experiments}
\label{sec:numerics}
%======================================================================

The numerical study retains two complementary manufactured benchmarks.
The first combines high oscillation with spatial localization on a square;
the second combines a reentrant-corner singularity with a separated
oscillatory packet on an L-shaped domain.  The exact solution is used only
for validation and never for marking, certification, or stopping.

At every exact checkpoint we report
\begin{equation}
\label{eq:numerical-certificate-components}
  \etaEq,
  \quad\underline\beta_\ell,
  \quad\rhoP,
  \quad\rhoA,
  \quad\epsP,
  \quad\epsA,
  \quad\betaInf,
  \quad\Uex,
\end{equation}
together with the true energy error and effectivity index
\begin{equation}
\label{eq:effectivity}
  E_{\mathrm{true}}:=\norm{u-U}_\kappa,
  \qquad
  I_{\mathrm{eff}}:=\frac{\Uex}{\norm{u-U}_\kappa}.
\end{equation}
A result is labelled \emph{certified} only when the geometric constants,
local algebraic solves, singular-value lower bound, and generalized
eigenvalue upper bounds are enclosed.  Otherwise the notation
$U_{\mathrm{ex}}^{\mathrm{pr}}$ is used.  Work measures include coarse and
fine degrees of freedom, oversampling level, number of cache batches,
number of refactorized patches, number of exact checkpoints, total flux
sweeps, and wall-clock time.

\subsection{Localized smooth oscillatory benchmark}
\label{sec:num-smooth}
%----------------------------------------------------------------------

Let
\[
  \Omega=(0,1)^2
\]
and choose the boundary decomposition
\begin{equation}
\label{eq:smooth-boundary-partition}
  \Gamma_D
  :=
  \{(x,0):0<x<1\}
  \cup
  \{(x,1):0<x<1\},
\end{equation}
\[
  \Gamma_N
  :=
  \{(0,y):0<y<1\},
  \qquad
  \Gamma_R
  :=
  \{(1,y):0<y<1\}.
\]
This is the same type of mixed homogeneous boundary condition as in
\eqref{eq:strong-problem}.

To generate a smooth but spatially localized oscillatory solution, set
\begin{equation}
\label{eq:smooth-manufactured}
  u(x,y)
  :=
  x^2(1-x)^2
  \sin(\pi y)
  \exp\!\left(
    -\alpha_{\mathrm{loc}}
    \bigl((x-x_\star)^2+(y-y_\star)^2\bigr)
  \right)
  e^{\ii\kappa x},
\end{equation}
where, for example,
\[
  (x_\star,y_\star)
  =
  \left(\frac34,\frac12\right),
  \qquad
  \alpha_{\mathrm{loc}}=80.
\]
The Gaussian factor localizes the dominant gradients near
$(x_\star,y_\star)$, while the factor $e^{\ii\kappa x}$ introduces the
oscillatory scale. This choice is intended to test whether the adaptive
procedure concentrates degrees of freedom in the region where they are
needed rather than refining the domain essentially uniformly.

The manufactured solution satisfies the homogeneous boundary conditions
of \eqref{eq:strong-problem}. Indeed,
\[
  u=0
  \qquad\text{on }\Gamma_D,
\]
and the factor $x^2(1-x)^2$ gives
\[
  \partial_xu=0
  \qquad\text{at }x=0,
\]
so that $\partial_\nu u=0$ on $\Gamma_N$. At $x=1$,
\[
  u=0,
  \qquad
  \partial_xu=0,
\]
and hence
\[
  \partial_\nu u-\ii\kappa u=0
  \qquad\text{on }\Gamma_R.
\]
The forcing is therefore defined solely by
\begin{equation}
\label{eq:smooth-forcing}
  f
  :=
  -\Delta u-\kappa^2u,
\end{equation}
with no additional boundary forcing.


We compare
\begin{enumerate}
\item the exact-certified adaptive LOD method of
      \cref{alg:exact-adaptive};
\item a fixed-$h$, fixed-oversampling LOD baseline;
\item uniform conforming $P_1$ finite elements;
\item adaptive conforming $P_1$ finite elements.
\end{enumerate}
For several wavenumbers, the initial coarse meshes are chosen so that the
same prescribed upper bound on $\kappa H$ is satisfied.  The experiment
reports error versus coarse and total local work, the effectivity index,
$\epsP$, $\epsA$, $\betaInf$, and the number of exact checkpoints.  The
coarse and fine meshes are plotted separately to show whether the local
Riesz marker concentrates coarse degrees of freedom near
$(x_\star,y_\star)$ and whether the total-defect marker refines only the
fine patches needed for certification.

The wavenumber study keeps $\alpha_{\mathrm{loc}}$ and the localization
center fixed and varies only the factor $e^{\ii\kappa x}$.  The main
robustness diagnostic is the certified prefactor
\begin{equation}
\label{eq:numerical-certified-prefactor}
  C_{\mathrm{cert}}
  :=c_W^{-1}
    +\frac{\epsA}{(1-\epsA)\betaInf}.
\end{equation}
Bounded behavior of $C_{\mathrm{cert}}$ over the tested range supports the
wavenumber-robust reliability theorem; it is not interpreted as a claim
that the total computational work is independent of $\kappa$.


\subsection{Low-regularity L-shaped benchmark}

\label{sec:num-lshape}

%----------------------------------------------------------------------

Let

\begin{equation}
\label{eq:lshape-domain}
  \Omega
  :=
  (-1,1)^2
  \setminus
  \bigl([0,1]\times[-1,0]\bigr),
\end{equation}

whose reentrant corner is located at the origin. Introduce polar
coordinates $(r,\theta)$ centered at the origin such that

\[
  0<\theta<\frac{3\pi}{2}
\]

locally inside $\Omega$.

The two boundary segments adjacent to the reentrant corner are assigned
Dirichlet conditions,

\begin{equation}
\label{eq:lshape-GammaD}
  \Gamma_D
  :=
  \{(x,0):0<x<1\}
  \cup
  \{(0,y):-1<y<0\}.
\end{equation}

We take

\begin{equation}
\label{eq:lshape-boundary-partition}
  \Gamma_N=\varnothing,
  \qquad
  \Gamma_R
  :=
  \partial\Omega\setminus\overline{\Gamma_D}.
\end{equation}

In particular, $|\Gamma_R|>0$, consistently with the standing assumption
of the model problem.

To retain the reentrant-corner singularity without introducing an
artificial radial transition layer, define the smooth boundary weight

\begin{equation}
\label{eq:lshape-boundary-weight}
  b_{\partial}(x,y)
  :=
  (1-x^2)^2(1-y^2)^2.
\end{equation}

The singular component is then chosen as

\begin{equation}
\label{eq:lshape-singularity}
  u_{\mathrm{sing}}(r,\theta)
  :=
  b_{\partial}(x,y)\,
  r^{2/3}
  \sin\!\left(\frac{2\theta}{3}\right).
\end{equation}

The factor

\[
  r^{2/3}\sin(2\theta/3)
\]

is the standard corner singularity associated with an interior angle
$3\pi/2$; see, e.g., \cite{BrennerScott2008}. It vanishes on the two rays
forming $\Gamma_D$ and has reduced regularity at the reentrant corner.
Since

\[
  b_{\partial}(x,y)=1+O(r^2)
  \qquad\text{as }r\to0,
\]

the multiplication by $b_{\partial}$ does not alter the leading-order
corner singularity. In contrast to a radial cut-off, the polynomial
weight does not introduce a distinguished transition annulus around the
origin. On the outer boundary it satisfies

\[
  b_{\partial}=0,
  \qquad
  \partial_\nu b_{\partial}=0
  \qquad\text{on }\Gamma_R.
\]

To introduce a spatially separated oscillatory feature, choose the
oscillation center

\begin{equation}
\label{eq:lshape-oscillation-center}
  (x_{\mathrm{osc}},y_{\mathrm{osc}})
  :=
  \left(-\frac12,\frac12\right)
\end{equation}

and define the Gaussian envelope

\begin{equation}
\label{eq:lshape-oscillation-envelope}
  G_{\mathrm{osc}}(x,y)
  :=
  \exp\!\left(
    -\alpha_{\mathrm{osc}}
    \bigl(
      (x-x_{\mathrm{osc}})^2
      +(y-y_{\mathrm{osc}})^2
    \bigr)
  \right).
\end{equation}

To enforce the homogeneous boundary conditions exactly, set

\begin{equation}
\label{eq:lshape-oscillation-amplitude}
  A_{\mathrm{osc}}(x,y)
  :=
  C_{\mathrm{osc}}\,
  xy\,
  b_{\partial}(x,y)\,
  G_{\mathrm{osc}}(x,y),
\end{equation}

where

\begin{equation}
\label{eq:lshape-oscillation-normalization}
  C_{\mathrm{osc}}
  :=
  \frac{1}{
    x_{\mathrm{osc}}y_{\mathrm{osc}}\,
    b_{\partial}(x_{\mathrm{osc}},y_{\mathrm{osc}})
  }.
\end{equation}

Thus

\[
  A_{\mathrm{osc}}
  (x_{\mathrm{osc}},y_{\mathrm{osc}})
  =1.
\]

The oscillatory component is

\begin{equation}
\label{eq:lshape-oscillation}
  u_{\mathrm{osc}}(x,y)
  :=
  \varepsilon_{\mathrm{osc}}\,
  A_{\mathrm{osc}}(x,y)
  e^{\ii\kappa(x-x_{\mathrm{osc}})}.
\end{equation}

Unless stated otherwise, we take

\begin{equation}
\label{eq:lshape-oscillation-parameters}
  \alpha_{\mathrm{osc}}=25,
  \qquad
  \varepsilon_{\mathrm{osc}}=0.25.
\end{equation}

The Gaussian envelope localizes the oscillatory scale around
$(x_{\mathrm{osc}},y_{\mathrm{osc}})$, which is separated from both the
reentrant corner and the outer boundary. The factor $xy$ makes
$u_{\mathrm{osc}}$ vanish on the two Dirichlet segments, while the
double zeros of $b_{\partial}$ imply

\[
  u_{\mathrm{osc}}=0,
  \qquad
  \partial_\nu u_{\mathrm{osc}}=0
  \qquad\text{on }\Gamma_R.
\]

The manufactured exact solution is

\begin{equation}
\label{eq:lshape-exact-solution}
  u
  :=
  u_{\mathrm{sing}}
  +
  u_{\mathrm{osc}}.
\end{equation}

For the singular component,
$r^{2/3}\sin(2\theta/3)$ vanishes on $\Gamma_D$, whereas
$u_{\mathrm{osc}}$ vanishes there because of the factor $xy$. Hence

\[
  u=0
  \qquad\text{on }\Gamma_D.
\]

Moreover, the properties of $b_{\partial}$ give

\[
  u=0,
  \qquad
  \partial_\nu u=0
  \qquad\text{on }\Gamma_R,
\]

and consequently

\[
  \partial_\nu u-\ii\kappa u=0
  \qquad\text{on }\Gamma_R.
\]

The volume forcing is generated by

\begin{equation}
\label{eq:lshape-forcing}
  f
  :=
  -\Delta u-\kappa^2u.
\end{equation}

For implementation, the forcing can be evaluated without differentiating
the corner singularity twice. Set

\[
  s(r,\theta)
  :=
  r^{2/3}\sin\!\left(\frac{2\theta}{3}\right).
\]

Since

\[
  \Delta s=0
  \qquad\text{for }r>0,
\]

the singular contribution to the forcing is

\begin{equation}
\label{eq:lshape-singular-forcing}
  f_{\mathrm{sing}}
  =
  -s\,\Delta b_{\partial}
  -2\nabla b_{\partial}\cdot\nabla s
  -\kappa^2 b_{\partial}s.
\end{equation}

Because
$\nabla b_{\partial}=O(r)$ near the origin and
$\nabla s=O(r^{-1/3})$, the right-hand side in
\eqref{eq:lshape-singular-forcing} is $O(r^{2/3})$ as $r\to0$ and admits
an $L^2$ extension through the corner. In particular, no forcing layer is
created at a prescribed positive distance from the singular point.

Similarly, using
\eqref{eq:lshape-oscillation}, the oscillatory contribution is

\begin{equation}
\label{eq:lshape-oscillatory-forcing}
  f_{\mathrm{osc}}
  =
  \varepsilon_{\mathrm{osc}}\,
  e^{\ii\kappa(x-x_{\mathrm{osc}})}
  \left(
    -\Delta A_{\mathrm{osc}}
    -2\ii\kappa\,
    \partial_x A_{\mathrm{osc}}
  \right),
\end{equation}

so that

\begin{equation}
\label{eq:lshape-forcing-split}
  f
  =
  f_{\mathrm{sing}}
  +
  f_{\mathrm{osc}}.
\end{equation}

No Neumann or Robin boundary forcing is introduced.


This benchmark tests whether the two-level adaptive mechanism separates
three tasks without a prescribed matching region: coarse approximation of
the corner singularity, coarse approximation of the localized oscillatory
packet, and fine-scale certification of the corrector space.  We compare
\begin{enumerate}
\item the exact-certified adaptive LOD method;
\item the same LOD method with fixed oversampling and a uniformly refined
      fine mesh;
\item adaptive conforming $P_1$ finite elements;
\item where useful, a uniform finite element baseline.
\end{enumerate}
The marked coarse elements should grade toward the reentrant corner and
also concentrate around $(x_{\mathrm{osc}},y_{\mathrm{osc}})$.  The
fine-defect indicator \eqref{eq:fine-defect-indicator} is inspected
separately; it need not reproduce the coarse mesh because its role is to
certify the continuous corrector defects rather than the primal
approximation error.

The reported histories include the exact error, $\Uex$, the effectivity
index, the localization diagnostic, the primal and adjoint dominant-mode
contributions, the distribution of coarse and fine cells, and the fraction
of local matrices reused within each cache batch.  Runs with different
$\ell$ and fine-mesh tolerances distinguish localization error from
fine-corrector discretization error.  No optimal-complexity conclusion is
drawn from the meshes alone.



%======================================================================
